% !TEX program = xelatex
% Odborna zprava podle dolozenych vysledku experimentu a E3 reportu.
% Doporucene umisteni v repozitari: docs/latex/main.tex
% Preklad: xelatex main.tex; pro grafy je dostupna zalozni vektorova varianta.
\documentclass[11pt,a4paper]{article}
\usepackage{fontspec}
\setmainfont{Liberation Serif}
\setsansfont{Liberation Sans}
\usepackage{polyglossia}
\setdefaultlanguage{czech}
\setotherlanguage{english}
\usepackage[a4paper,top=18mm,bottom=18mm,left=23mm,right=23mm]{geometry}
\usepackage{microtype}
\usepackage{amsmath,amssymb,mathtools}
\usepackage{graphicx,booktabs,tabularx,multirow,array}
\usepackage{caption}
\captionsetup{font=small,labelfont=bf,skip=4pt}
\usepackage{float}
\usepackage{needspace}
\usepackage{enumitem}
\usepackage{titlesec}
\usepackage{tikz}
\usepackage{pgfplots}
\pgfplotsset{compat=1.18}
\usepackage[unicode,colorlinks=true,linkcolor=black,urlcolor=blue!50!black,citecolor=black]{hyperref}
\usepackage{url}
\setlength{\parindent}{0pt}
\setlength{\parskip}{4.8pt plus 1pt minus 0.8pt}
\setlength{\textfloatsep}{8pt plus 2pt minus 2pt}
\setlength{\intextsep}{7pt plus 2pt minus 2pt}
\setlength{\tabcolsep}{5.2pt}
\renewcommand{\arraystretch}{1.12}
\titleformat{\section}{\Large\bfseries\sffamily}{\thesection.}{0.5em}{}
\titleformat{\subsection}{\normalsize\bfseries\sffamily}{\thesubsection}{0.55em}{}
\titlespacing*{\section}{0pt}{13pt}{5pt}
\titlespacing*{\subsection}{0pt}{10pt}{3pt}
\setlist[itemize]{topsep=2pt,itemsep=1pt,parsep=0pt,leftmargin=17pt}
\newcommand{\ap}{\mathrm{AP}}
\newcommand{\roc}{\mathrm{ROC\text{-}AUC}}
\newcommand{\dataset}{\texttt{saadtaleb/tess-lightcurves-planets}}
\newcommand{\code}[1]{\texttt{#1}}

\begin{document}
\begin{center}
  {\Large\bfseries\sffamily Klasifikace tranzitních signálů TESS pomocí hlubokého učení}\par
  \vspace{2.5mm}
  {\normalsize Srovnání vlastní 1D CNN a předtrénovaného modelu Mantis-8M}\par
  \vspace{2.0mm}
  {\small Semestrální projekt: Úvod do deep learningu \quad|\quad listopad 2026}\par
  {\small Autor: Vladimír Pospíšil}\par
\end{center}
\vspace{-1.5mm}
\begin{abstract}
\noindent
Práce zkoumá binární klasifikaci již katalogově označených tranzitních signálů z fotometrických měření družice TESS. Porovnává logistickou regresi, vlastní dvouvětvovou jednorozměrnou konvoluční síť (CNN) a předtrénovaný encoder Mantis-8M se samostatně učenou klasifikační hlavou. Z~13\,254 katalogových záznamů bylo vybráno 1\,918 použitelných signálů se štítkem potvrzená planeta nebo falešně pozitivní signál. Hvězdy byly rozděleny podle identifikátoru TIC tak, aby se nevyskytovaly v~různých množinách. Na jednorázově vyhodnoceném nezávislém testu 290 signálů dosáhly modely v metrice \emph{Average Precision} (AP) 0,7981, 0,8851\,$\pm$\,0,0049 a 0,9009\,$\pm$\,0,0024. Doplňková pětinásobná skupinová CV nad vývojovými daty podpořila vyšší průměrnou AP modelu Mantis; přínos masky pokrytí se nepotvrdil. Výsledky podporují použitelnost přenosu reprezentací z předtrénovaného modelu, ale nevypovídají o spolehlivosti potvrzování dosud neznámých exoplanet.
\end{abstract}

\section{Úvod a cíl práce}
Exoplanety lze zjišťovat mimo jiné pomocí tranzitní fotometrie: při přechodu tělesa před diskem mateřské hvězdy dojde k malému periodickému poklesu jejího pozorovaného jasu. Družice \emph{Transiting Exoplanet Survey Satellite} (TESS) poskytuje rozsáhlé světelné křivky hvězd vhodné pro vyhledávání takových jevů \cite{ricker2015}; soubory FITS jsou zpřístupněny v archivu MAST \cite{mast}. Pozorovaný pokles však může mít také neplanetární příčinu, například zákrytovou dvojhvězdu nebo instrumentální systematiku.

Cílem práce není hledat nové periodické události v~neznámé časové řadě, ale \textbf{rozlišovat již identifikované signály} se dvěma katalogovými štítky: potvrzená planeta a falešně pozitivní signál. Výzkumná otázka zní, zda model založený na předtrénovaných reprezentacích dosáhne při malém počtu štítkovaných příkladů lepšího výsledku než vlastní CNN a klasický baseline. Doplňující experimenty zkoumají význam lokálního a globálního pohledu i explicitní informace o pokrytí jednotlivých fázových intervalů.

\section{Data a předzpracování}
\subsection{Zdroj a výběr signálů}
Použit byl soubor \dataset\ publikovaný na Hugging Face \cite{dataset}. Datový soubor přebírá katalogové štítky mimo jiné ze služeb ExoFOP-TESS a NASA Exoplanet Archive \cite{exofop,nea}; jeho veřejná karta uvádí licenci CC0 1.0. Štítky jsou katalogové, nikoli výsledky vlastní ruční anotace. Metadata zahrnují 13\,254 signálů: 9\,895 kandidátů, 1\,312 potvrzených planet, 666 falešně pozitivních a 1\,381 záznamů s~neznámým štítkem. Neověřené kandidáty nelze bez dalšího považovat za pozitivní příklady, a~proto byly ze supervidované binární úlohy vyloučeny.

Výchozích 1\,978 známě označených signálů prošlo postupným filtrováním. Nejdříve 39 příkladů nesplnilo souhrnnou podmínku \code{preliminary\_ready}, zahrnující kontroly metadat, efemeridy a dostupnosti použitelného FITS; dalších 7 neprošlo kontrolou tranzitního pokrytí, 1 měl délku tranzitu alespoň čtvrtinu periody a 13 nemělo dostatečné pokrytí fázových binů. Vzniklo \textbf{1\,918 signálů}: 1\,286 potvrzených planet a 632 falešně pozitivních případů. Rozpad prvních 39 vyřazených příkladů podle jednotlivých příčin nebyl dosud samostatně kvantifikován.

\begin{table}[H]
\centering\small
\caption{Velikost množin po skupinovém dělení podle TIC.}
\label{tab:splits}
\begin{tabular}{lrr}
\toprule
Množina & Počet signálů & Unikátní TIC \\
\midrule
Trénovací & 1\,335 & 1\,252 \\
Validační & 293 & 268 \\
Testovací & 290 & 269 \\
\midrule
Celkem & \textbf{1\,918} & \textbf{1\,789} \\
\bottomrule
\end{tabular}
\end{table}

Dělení pomocí TIC skupin odpovídá situaci, kdy jedna hvězda může mít více signálů a více sektorových záznamů. Překryv hvězd by mohl vést k~úniku hvězdně specifických informací mezi trénováním a evaluací. Audit ověřil nepřekryvnost TIC i indexů a úplné přiřazení všech signálů.

\subsection{Transformace světelných křivek}
Pro každý signál byly ze souborů FITS načteny čas, fotometrický tok a podle dostupnosti příznaky kvality. Nepoužitelná měření byla odstraněna, světelný tok byl normalizován vůči referenční hladině mimo tranzit a čas převeden do společné soustavy vzhledem k~epoše $t_0$ a periodě $P$. Fázová souřadnice vzniká ze vztahu
\begin{equation}
 \phi(t)=\frac{\big[(t-t_0+P/2)\bmod P\big]-P/2}{P}.
 \label{eq:phase}
\end{equation}
Lokální pohled pokrývá interval $\pm2{,}5$ násobku délky tranzitu a obsahuje $101$ binů; globální pohled zahrnuje celou fázi $[-0{,}5,0{,}5)$ v~$201$ binech. V~každém binu se počítá medián relativního toku, pokud jsou k~dispozici alespoň tři měření. Druhým kanálem je binární maska pozorovanosti. Požadováno je alespoň 45\,\% pokrytých lokálních, 30\,\% globálních a pět centrálních lokálních binů. Vstupy mají rozměry $101\times2$ a $201\times2$.

\section{Modely a matematická formulace}
Uvažujme dataset $\mathcal D=\{(x_i,y_i)\}_{i=1}^{n}$, kde $y_i\in\{0,1\}$ ($1$ = potvrzená planeta, $0$ = falešně pozitivní signál) a $x_i=(x_i^{\mathrm{loc}},x_i^{\mathrm{glob}})$. Model $f_{\theta}$ vytváří \emph{logit} $z_i=f_{\theta}(x_i)\in\mathbb R$, který se převádí sigmoidou na skóre $p_i=\sigma(z_i)=1/(1+e^{-z_i})$. Pro neuronové modely byla použita vážená binární křížová entropie nad logity:
\begin{equation}
\mathcal L(\theta)=-\frac1n\sum_{i=1}^n\!\left[w_+ y_i\log\sigma(z_i)+(1-y_i)\log\bigl(1-\sigma(z_i)\bigr)\right],
\qquad w_+=\frac{n_-}{n_+},
\label{eq:loss}
\end{equation}
kde $n_+$ a $n_-$ jsou počty pozitivních a negativních trénovacích příkladů. V~implementaci je použita numericky stabilní \code{BCEWithLogitsLoss} s~\code{pos\_weight}. Minimalizaci provádí AdamW. Výstupní skóre nejsou automaticky kalibrované pravděpodobnosti astronomického potvrzení.

\subsection{Baseline: logistická regrese}
Jako srovnávací klasický model slouží logistická regrese nad zploštěnými lokálními a globálními vstupy (celkem 604 příznaků). Standardizace je odvozena pouze z~trénovacích dat. Použita byla logistická regrese s~$L_2$ regularizací ($C=0{,}1$), váhami tříd \code{balanced} a solverem \code{liblinear}. Baseline dovoluje posoudit, zda neuronové reprezentace přidávají užitek oproti lineární rozhodovací hranici nad stejnými pohledy.

\subsection{Model A: vlastní dvouvětvová 1D CNN}
CNN byla navržena a trénována od náhodné inicializace. Každý pohled zpracovává samostatná větev s~konvolucemi $2\!\rightarrow\!32$ (jádro 7), $32\!\rightarrow\!64$ (jádro 5) a $64\!\rightarrow\!64$ (jádro 3), následovanými normalizací, ReLU a dvěma operacemi MaxPool po prvních vrstvách. Adaptivní průměrový a maximový pooling vytvoří v~každé větvi 128 příznaků. Jejich konkatenace dává vektor délky 256, který prochází vrstvou $256\rightarrow128$, ReLU, dropoutem a lineárním výstupem $128\rightarrow1$. Celkový počet parametrů je \textbf{79\,937}.

Před vstupem se relativní tok transformuje pomocí $\operatorname{arsinh}(1000f/5)$; maska pokrytí se zachovává. Dvě větve umožňují kombinovat jemný tvar tranzitu v okolí minima s~kontextem celé orbitální fáze.

\begin{figure}[H]
\centering
\begin{tikzpicture}[
 box/.style={draw=black!55,rounded corners=3pt,align=center,minimum height=8mm,fill=black!3,font=\small},
 arr/.style={-stealth,thick,draw=black!65}]
\node[box,minimum width=43mm] (l) at (0,2.5) {Lokální pohled\\$101\times2$};
\node[box,minimum width=43mm] (g) at (8.5,2.5) {Globální pohled\\$201\times2$};
\node[box,minimum width=43mm] (le) at (0,1) {3$\times$ Conv1D + BN + ReLU\\avg/max pooling: 128};
\node[box,minimum width=43mm] (ge) at (8.5,1) {3$\times$ Conv1D + BN + ReLU\\avg/max pooling: 128};
\node[box,minimum width=50mm] (cat) at (4.25,-0.4) {Konkatenace: 256 příznaků};
\node[box,minimum width=50mm] (head) at (4.25,-1.55) {MLP: 256 $\to$ 128 $\to$ 1};
\node[box,minimum width=50mm,fill=black!8] (out) at (4.25,-2.7) {Logit $z$ a skóre $\sigma(z)$};
\draw[arr] (l) -- (le); \draw[arr] (g) -- (ge);
\draw[arr] (le) -- (cat); \draw[arr] (ge) -- (cat);
\draw[arr] (cat) -- (head); \draw[arr] (head) -- (out);
\end{tikzpicture}
\caption{Schéma dvouvětvové CNN; obě větve mají vlastní váhy.}
\label{fig:architecture}
\end{figure}

\subsection{Model B: Mantis-8M a klasifikační hlava}
Model B používá existující předtrénovaný časověřadový encoder \emph{Mantis-8M} ze starší varianty MantisV1, načtený z~checkpointu \code{paris-noah/Mantis-8M} \cite{mantisoriginal,mantis2026,mantischeckpoint}. Citovaná práce ICML 2026 popisuje rodinu Mantis; \textbf{nepoužili jsme novější checkpoint MantisV2}. Encoder má \textbf{8\,103\,936} parametrů a zůstává zmrazený. Lokální a globální tok se transformují pomocí $\operatorname{arsinh}(f/5)$ a samostatně interpolují na délku 512. Dvě 512rozměrné reprezentace se spojí do vektoru délky 1\,024. Na něm se trénuje MLP hlava $1024\rightarrow128\rightarrow1$ s~ReLU a dropoutem; hlava má \textbf{131\,329} trénovaných parametrů. Standardizace embeddingů se odhaduje výhradně na trénovací množině. Pro zmrazený encoder byl použit výstup z třetí (index 2) transformerové vrstvy s nastavením \code{output\_token=combined}, v souladu s dokumentací Mantis.

Srovnání A a B není srovnáním shodného objemu trénovacích dat a výpočetního rozpočtu: předtrénovaný encoder již využívá reprezentace naučené v~předchozí fázi.

\section{Experimentální protokol}
Hlavní metrikou pro výběr konfigurace byla \emph{Average Precision} (AP), tedy souhrn precision při změnách recall. U~nevyvážených tříd je vhodné hodnotit také PR chování, nikoli pouze ROC \cite{saito2015}. Doplňkově byly měřeny ROC-AUC, vyvážená přesnost a F1 při pevném prahu 0,5. Testovací množina byla před výběrem architektur, hyperparametrů a checkpointů uzavřena; finální evaluace všech zmrazených checkpointů proběhla v jediném testovacím průchodu.

Pro popis klasifikace při daném prahu jsou použity počty $TP$, $TN$, $FP$ a $FN$; definujeme $\mathrm{Precision}=TP/(TP+FP)$, $\mathrm{Recall}=TP/(TP+FN)$, $F_1=2PR/(P+R)$ a $\mathrm{BA}=(TPR+TNR)/2$. Základní orientací pro AP je prevalence pozitivní třídy, tedy podíl planet na dané evaluační množině. Během ladění jsme však nesměli využít testovací štítky ani výsledné testovací AP.

Pro CNN proběhl grid search $3\times3\times3=27$ konfigurací: learning rate $\{0{,}0003;0{,}001;0{,}003\}$, dropout $\{0{,}1;0{,}3;0{,}5\}$ a batch size $\{32;64;128\}$. Nejlepší konfigurace podle validační AP měla $\eta=0{,}001$, dropout 0,3, batch 32, weight decay $10^{-4}$, maximum 35 epoch a trpělivost early stopping 7 epoch. Pro Mantis proběhl grid devíti kombinací skryté šířky $\{64,128,256\}$ a dropoutu $\{0{,}1;0{,}3;0{,}5\}$; vybrána byla šířka 128 a dropout 0,5. Pro hlavu platí learning rate $10^{-3}$, weight decay $10^{-3}$, batch 64, nejvýše 80 epoch a patience 12.

Vybrané konfigurace byly zopakovány se seedy $42$, $123$ a $2026$. Všechny běhy sdílely jediný pevný TIC-disjunktní split. Směrodatná odchylka proto charakterizuje citlivost na inicializaci, nikoli nejistotu při nezávislém výběru jiných hvězd. Seed 42 se současně účastnil výběru hyperparametrů.

\subsection{Dva cílené experimenty}
\textbf{E1 -- příspěvek obou fázových pohledů.} Hypotéza: kombinace lokální a globální části křivky zlepší validační AP. Zvolena byla \emph{ablace na úrovni příznaků}: výstup nepoužité větve CNN byl při trénování i validaci nahrazen nulami, aby její konvoluční bias a BatchNorm nevytvářely nenulové konstantní příznaky. Data, rozdělení podle TIC, optimalizátor, hyperparametry a seedy zůstaly stejné. Neaktivní větev však nepřispívá gradientem, proto jednotlivé varianty nemají stejnou \emph{efektivní} kapacitu. Druhý, samostatně uložený pokus s nulováním vstupu (viz \code{docs/model\_a\_experiment\_e1\_views.md}) je pomocnou kontrolou, nikoli zdrojem následující tabulky.

\begin{table}[H]
\centering\small
\caption{E1: validační metriky (průměr $\pm$ výběrová směrodatná odchylka, tři seedy); ablace výstupních příznaků větví.}
\label{tab:e1}
\begin{tabular}{lcc}
\toprule
Použité pohledy & AP & F1 při prahu 0,5 \\
\midrule
Oba pohledy & $\mathbf{0{,}9219\pm0{,}0127}$ & $0{,}8131\pm0{,}0300$ \\
Pouze lokální & $0{,}8994\pm0{,}0027$ & $0{,}7368\pm0{,}0512$ \\
Pouze globální & $0{,}8933\pm0{,}0041$ & $0{,}6819\pm0{,}0389$ \\
\bottomrule
\end{tabular}
\end{table}

Společné použití obou pohledů přidalo na validaci $+0{,}0225$ AP oproti lokálnímu a $+0{,}0286$ oproti globálnímu pohledu. Všechny tři párované seedy měly kladný rozdíl, což podporuje hypotézu pro tento výběr dat; nejde o nezávislé potvrzení na testu.

\textbf{E2 -- informace o pokrytí.} Hypotéza: explicitní kanál masky zlepší AP. V ablační variantě se maskový kanál nahradil nulami během trénování i validace, zatímco flux, architektura, split a hyperparametry zůstaly stejné. Průměrná validační AP se změnila z~$0{,}9205\pm0{,}0114$ bez masky na $0{,}9219\pm0{,}0127$ s maskou ($\Delta=+0{,}0014$). Zlepšení je malé a jeden ze tří seedů se nepatrně zhoršil. Nulové hodnoty fluxových binů navíc mohou samy zčásti kódovat chybějící pozorování.

\begin{table}[H]
\centering\small
\caption{E2: validační výsledky ablace masky (stejný TIC split, tři seedy).}
\label{tab:e2}
\begin{tabular}{lcc}
\toprule
Metrika & S maskou pokrytí & Bez masky \\
\midrule
AP & $0{,}9219\pm0{,}0127$ & $0{,}9205\pm0{,}0114$ \\
ROC-AUC & $0{,}8597\pm0{,}0054$ & $0{,}8532\pm0{,}0193$ \\
Vyvážená přesnost & $0{,}7685\pm0{,}0428$ & $0{,}7832\pm0{,}0296$ \\
F1 & $0{,}8131\pm0{,}0300$ & $0{,}7994\pm0{,}0493$ \\
\bottomrule
\end{tabular}
\end{table}

Oba experimenty používají validační množinu, na níž proběhl i výběr hyperparametrů; výsledky proto nelze chápat jako další nezávislý test. Jako techniky nad rámec základního trénování jsou použity dvoupohledová fázová reprezentace, ablace vstupů a přenos reprezentací z předtrénovaného encoderu.

\section{Výsledky}
\begin{table}[H]
\centering\small
\caption{Jednorázový test na 290 signálech; neurální modely: průměr $\pm$ výběrová směrodatná odchylka pro tři seedy.}
\label{tab:results}
\resizebox{\linewidth}{!}{%
\begin{tabular}{lcccc}
\toprule
Model & AP & ROC-AUC & Vyvážená přesnost & F1 \\
\midrule
Logistická regrese & 0,7981 & 0,6979 & 0,6541 & 0,6189 \\
CNN (A) & $0{,}8851\pm0{,}0049$ & $0{,}8225\pm0{,}0014$ & $0{,}7308\pm0{,}0107$ & $0{,}7772\pm0{,}0679$ \\
Mantis + MLP (B) & $\mathbf{0{,}9009\pm0{,}0024}$ & $\mathbf{0{,}8434\pm0{,}0018}$ & $\mathbf{0{,}7539\pm0{,}0144}$ & $\mathbf{0{,}8119\pm0{,}0106}$ \\
\bottomrule
\end{tabular}}
\end{table}

Testovací množina obsahovala 195 potvrzených planet a 95 falešně pozitivních signálů. Podíl pozitivní třídy $195/290\doteq0{,}672$ je nutné zohlednit při interpretaci AP. Mantis dosáhl proti CNN rozdílu průměrné AP $+0{,}0158$ a proti logistické regresi $+0{,}1028$; z jednoho testovacího splitu však nelze dovozovat statisticky významnou převahu na populaci všech kandidátů TESS.

\begin{figure}[H]
\centering
% Pokud existuje puvodni graf v repozitari, pouzijeme ho. Jinak samostatny vektorovy graf.
\IfFileExists{../figures/final_test_ap_comparison.pdf}{%
\includegraphics[width=0.70\linewidth]{../figures/final_test_ap_comparison.pdf}%
}{%
\begin{tikzpicture}
\begin{axis}[
 width=0.83\linewidth,height=45mm,ymin=0,ymax=1,
 symbolic x coords={Baseline,CNN,Mantis},xtick=data,
 ylabel={Average Precision},ymajorgrids=true,
 grid style={black!10},bar width=19pt,
 nodes near coords,nodes near coords align={vertical},
 every node near coord/.append style={font=\scriptsize},
 tick label style={font=\small},ylabel style={font=\small},
 ]
\addplot+[ybar,fill=black!60,draw=black!65] coordinates {(Baseline,0.7981) (CNN,0.8851) (Mantis,0.9009)};
\addplot+[no marks,dashed,black!60] coordinates {(Baseline,0.6724) (Mantis,0.6724)};
\end{axis}
\end{tikzpicture}
}
\caption{Výsledná Average Precision na nezávislé testovací množině. Přerušovaná čára vyznačuje prevalenci pozitivní třídy (0,672).}
\label{fig:ap}
\end{figure}

\begin{table}[H]
\centering\small
\caption{Porovnání validační a testovací AP po výběru konfigurací.}
\label{tab:valtest}
\begin{tabular}{lrrr}
\toprule
Model & Validace AP & Test AP & Rozdíl (test--val.) \\
\midrule
CNN & 0,9219 & 0,8851 & $-0{,}0368$ \\
Mantis + MLP & 0,9420 & 0,9009 & $-0{,}0411$ \\
\bottomrule
\end{tabular}
\end{table}

\begin{table}[H]
\centering\small
\caption{Jednotlivé běhy na testovacím souboru.}
\label{tab:seeds}
\begin{tabular}{lcrr}
\toprule
Model & Seed & AP & F1 \\
\midrule
Baseline & -- & 0,7981 & 0,6189 \\
\midrule
CNN & 42 & 0,8837 & 0,7727 \\
CNN & 123 & 0,8809 & 0,8473 \\
CNN & 2026 & 0,8905 & 0,7117 \\
\midrule
Mantis & 42 & 0,8984 & 0,8000 \\
Mantis & 123 & 0,9030 & 0,8204 \\
Mantis & 2026 & 0,9014 & 0,8152 \\
\bottomrule
\end{tabular}
\end{table}

Na testu se AP u CNN mezi seedy pohybovala od 0,8809 do 0,8905, u Mantis od 0,8984 do 0,9030. Naproti tomu F1 CNN vykazovala větší variabilitu mezi inicializacemi ($0{,}7772\pm0{,}0679$) než Mantis ($0{,}8119\pm0{,}0106$), přestože AP byla relativně stabilní. To ukazuje citlivost pevného rozhodovacího prahu na rozložení výstupních skóre různých inicializací.

\begin{figure}[H]
\centering
\IfFileExists{../figures/final_test_f1_by_seed.pdf}{%
\includegraphics[width=0.80\linewidth]{../figures/final_test_f1_by_seed.pdf}%
}{%
\begin{tikzpicture}
\begin{axis}[
 width=0.82\linewidth,height=49mm,ymin=0.55,ymax=0.9,
 symbolic x coords={42,123,2026},xtick=data,
 xlabel={Náhodný seed},ylabel={F1-score},
 ymajorgrids=true,grid style={black!12},
 legend pos=south east,
 tick label style={font=\small},
 ]
\addplot+[mark=*,thick] coordinates {(42,0.7727) (123,0.8473) (2026,0.7117)};
\addlegendentry{CNN}
\addplot+[mark=square*,thick,dashed] coordinates {(42,0.8000) (123,0.8204) (2026,0.8152)};
\addlegendentry{Mantis}
\end{axis}
\end{tikzpicture}
}
\caption{Proměnlivost F1 mezi inicializacemi při pevném rozhodovacím prahu 0,5.}
\label{fig:f1}
\end{figure}

Uloženy byly rovněž počty záměn. Například Mantis při seedu 123 správně klasifikoval 153 potvrzených a 70 falešně pozitivních signálů; 42 potvrzených označil jako negativní a 25 falešně pozitivních jako pozitivní. Při interpretaci nesmí být opomenuto, že katalogové štítky jsou kontextově podmíněné a nemusí odpovídat definitivní fyzikální charakterizaci všech jevů.


\Needspace{8\baselineskip}
\subsection{Doplňkové ověření: pětinásobná skupinová cross-validation (E3)}
Pro ověření citlivosti výsledků na složení vývojových hvězd byla po dokončení původního experimentu provedena doplňková \emph{5-fold} cross-validation nad 1\,628 původními trénovacími a validačními signály. Původní nezávislý test 290 signálů zůstal oddělený. \code{StratifiedGroupKFold} seskupoval jednotlivé signály podle TIC, takže tatáž hvězda nemohla vstoupit současně do trénovací a hodnocené části foldu. Pět validačních foldů obsahovalo 317--334 signálů; logistická regrese byla vyhodnocena jednou na fold, CNN, Mantis i ablační varianty pro tři seedy; u neuronových modelů se epocha vybírala na vnitřní TIC-validaci a následně se model znovu učil na celé trénovací části daného vnějšího foldu. Standardizace byla v každém foldu odvozena pouze z jeho trénovacích dat.

\begin{table}[H]
\centering\footnotesize
\caption{E3: průměr a výběrová SD z pěti vnějších foldů (neuronové modely: nejprve průměr tří seedů v každém foldu); BA značí vyváženou přesnost.}
\label{tab:e3models}
{\setlength{\tabcolsep}{3pt}
\begin{tabular}{lcccc}
\toprule
Model & AP & ROC-AUC & BA & F1 \\
\midrule
Logistická regrese & $0{,}8016\pm0{,}0247$ & $0{,}7114\pm0{,}0394$ & $0{,}6726\pm0{,}0361$ & $0{,}6756\pm0{,}0612$ \\
CNN & $0{,}9000\pm0{,}0151$ & $0{,}8321\pm0{,}0240$ & $0{,}7414\pm0{,}0317$ & $0{,}7784\pm0{,}0507$ \\
Mantis-8M + MLP & $\mathbf{0{,}9289\pm0{,}0071}$ & $0{,}8797\pm0{,}0109$ & $0{,}7947\pm0{,}0232$ & $0{,}8452\pm0{,}0158$ \\
\bottomrule
\end{tabular}}
\end{table}

V E3 měl Mantis vyšší průměrnou AP než CNN ve všech pěti foldech. Párová ablace E1 podpořila přínos obou pohledů: CNN s oběma dosáhla AP $0{,}9000$, pouze s lokálním $0{,}8855$ a pouze s globálním $0{,}8729$; rozdíly činily $+0{,}0145\pm0{,}0148$, respektive $+0{,}0271\pm0{,}0170$ (kladné ve 12/15, respektive 13/15 párech; SD mezi pěti foldy). U E2 se však malý původní pozitivní efekt masky nepotvrdil: AP s maskou $0{,}9000$, bez masky $0{,}9015$, párový rozdíl $-0{,}0014\pm0{,}0040$ ve prospěch varianty bez masky (počítaný z nezaokrouhlených hodnot; SD mezi pěti foldy). Maska měla vyšší AP jen v 8 z 15 párovaných běhů. E3 je dodatečná analýza robustnosti, nikoli nový nezávislý test nebo nezaujatý odhad po výběru hyperparametrů; ty byly před CV vybrány na části vývojových dat. Uváděná SD shrnuje pět průměrů foldů, ne 15 nezávislých datových splitů.

\section{Diskuse, omezení a náklady}
Silnější výkon Mantis může být dán vhodnější reprezentací časových řad, ale také výhodou externího předtrénování; na základě této studie nelze obě vysvětlení oddělit. Malé zastoupení negativních příkladů, selekce pouze dvou katalogových štítků a vysoká pozitivní prevalence omezují přenositelnost výsledků na nepotvrzené kandidáty. Původní test obsahuje jediný pevný split a tři seedy, takže neposkytuje dostatečný podklad k testu statistické významnosti malého rozdílu AP. Doplňková CV (E3) omezuje závislost závěrů na jednom vývojovém rozdělení, ale nepředstavuje nový nezávislý test.

Analýza validačních chyb upozornila na opakovaná falešná pozitiva u několika systémů označených jako zákrytové dvojhvězdy a na falešná negativa i u signálů se zřetelným tranzitem. Smysluplná následná práce by měla zahrnovat kontrolu nezávislých astronomických katalogů, systematické hodnocení nejistoty a případnou externí sadu hvězd. Skóre modelu není astronomickým potvrzením planety.

Výpočetní prostředí: Windows 11, Python 3.10.6, NVIDIA GeForce RTX 4060 Laptop GPU a CUDA, PyTorch 2.14.1+cu130. CNN má 79\,937 parametrů (všechny trénované); Mantis má 8\,103\,936 zmrazených parametrů encoderu a 131\,329 trénovaných parametrů hlavy, celkem 8\,235\,265. Uloženy jsou trénovací/validační křivky CNN s loss a AP a řady metrik během trénování; ukázková trénovací křivka je dostupná v \code{docs/figures/}, zatímco ostatní křivky a checkpointy jednotlivých běhů zůstávají lokálně ve \code{full\_dataset/}.

Trénování finální CNN trvalo podle seedu 5,26 / 8,89 / 4,35~s. Pro \emph{původní} netuněnou hlavu Mantis jsou zdokumentovány časy 1,14 / 1,15 / 0,94~s (bez extrakce embeddingů); pro \emph{vybranou} doladěnou hlavu nejsou v commitnutém shrnutí časy uvedeny, takže je nelze věrohodně porovnat. Jejich přesné hodnoty lze dodatečně odečíst z místních \code{metrics.json}, pokud se dochovaly. Výpočetní čas externího předtrénování Mantis není znám.

Při testu 290 signálů trvala extrakce dvou pohledů Mantis přibližně 0,503~s (batch 16), inference MLP hlavy asi 0,002~s na seed. Logistická regrese predikovala za 0,002~s. CNN měla naměřeno 0,247 / 0,018 / 0,006~s (batch 64), což ilustruje výrazný vliv zahřátí a dalších podmínek měření; nepovažujeme je za přesně srovnatelný benchmark. Uvedené časy nezahrnují načtení modelů.
Reprodukovatelnost podporují zdrojové skripty, pevná náhodná semena, konfigurační soubory, kontrola SHA-256 vstupů a zaznamenané metriky v~GitHub repozitáři. Rozsáhlé FITS, odvozené křivky trénování a checkpointy nejsou uloženy v~Gitu. Kompletní nezávislý rerun od čistého klonu nebyl proveden; popsaný postup a verze závislostí proto představují dokumentované, nikoli experimentálně potvrzené záruky úplné reprodukce.

\section{Závěr}
Studie splnila cíl porovnat jednoduchý baseline, vlastní 1D CNN a adaptovaný předtrénovaný encoder na konzistentně připravené binární úloze TESS. Všechny modely byly posouzeny na společném TIC-disjunktním testu. Nejvyšší průměrnou AP $0{,}9009\pm0{,}0024$ získal Mantis-8M s~MLP hlavou, vlastní CNN dosáhla $0{,}8851\pm0{,}0049$ a logistická regrese 0,7981. Experiment E1 a doplňková 5-fold CV podpořily použití obou fázových pohledů; E2 vykázal na původní validaci malý pozitivní efekt masky, který se v CV nepotvrdil. Obě zjištění zdůrazňují význam kontrolovaných ablačních studií. Navazujícím samostatným směrem je inference na nepotvrzených kandidátech, nikoli jejich automatické prohlašování za planety.

\paragraph{Využití generativní AI.} Generativní AI byla použita při návrhu experimentů, formulování a kontrole kódu, ladění, dokumentaci a sepisování první verze zprávy. Správnost výstupů byla ověřována spuštěním skriptů, kontrolou katalogových štítků a TIC, auditem uložených artefaktů, validačními běhy a jednorázovým testem. Za interpretaci a konečnou odbornou správnost odpovídá autor.

\vspace{-2pt}
\begin{thebibliography}{10}
\small
\setlength{\itemsep}{1pt}
\bibitem{ricker2015} G.~R. Ricker et al.: Transiting Exoplanet Survey Satellite (TESS). \emph{Journal of Astronomical Telescopes, Instruments, and Systems} 1(1), 014003 (2015). \href{https://doi.org/10.1117/1.JATIS.1.1.014003}{doi:10.1117/1.JATIS.1.1.014003}.
\bibitem{dataset} Saadtaleb: \emph{TESS Light Curves for Planet Detection}. Datový soubor, Hugging Face (2026). \url{https://huggingface.co/datasets/saadtaleb/tess-lightcurves-planets}.
\bibitem{mast} STScI: \emph{TESS at MAST}. Archivační služba. \url{https://archive.stsci.edu/missions-and-data/tess}.
\bibitem{exofop} IPAC: \emph{Exoplanet Follow-up Observing Program -- TESS}. \url{https://exofop.ipac.caltech.edu/tess/}.
\bibitem{nea} NASA Exoplanet Science Institute: \emph{NASA Exoplanet Archive}. \url{https://exoplanetarchive.ipac.caltech.edu/}.
\bibitem{mantis2026} V.~Feofanov et al.: Mantis: Lightweight Foundation Model for Time Series Classification. \emph{Proceedings of the 43rd International Conference on Machine Learning}, PMLR 306, 30692--30719 (2026). \url{https://proceedings.mlr.press/v306/feofanov26a.html}.
\bibitem{mantisoriginal} V.~Feofanov et al.: Mantis: Lightweight Calibrated Foundation Model for User-Friendly Time Series Classification. \emph{arXiv:2502.15637} (2025). \url{https://arxiv.org/abs/2502.15637}.
\bibitem{mantischeckpoint} Huawei Paris Noah's Ark Lab: \emph{Mantis-8M} (MantisV1), modelová karta a checkpoint. Hugging Face. \url{https://huggingface.co/paris-noah/Mantis-8M}.
\bibitem{saito2015} T.~Saito, M.~Rehmsmeier: The Precision-Recall Plot Is More Informative than the ROC Plot When Evaluating Binary Classifiers on Imbalanced Datasets. \emph{PLOS ONE} 10(3), e0118432 (2015). \href{https://doi.org/10.1371/journal.pone.0118432}{doi:10.1371/journal.pone.0118432}.
\end{thebibliography}
\end{document}
